\chapter{Revisão de Literatura}
\label{cap:RefTeorico}

Este capítulo busca familiarizar o leitor com conceitos e termos técnicos a serem utilizados no trabalho, fornecendo o conhecimento mínimo que o leitor precisa ter para entender o que será estudado mais adiante. São apresentados conceitos já existentes na literatura, sempre realizando as devidas citações. Desse modo, o referencial teórico possibilita a contextualização do trabalho a ser elaborado e uniformiza o conhecimento necessário para o entendimento completo do TCC.

O autor deve mencionar a literatura que serviu de base ao desenvolvimento da pesquisa, citando apenas as referências pertinentes ao assunto do trabalho e limitando-se, dentro do possível, às mais atualizadas. Não deve ser uma simples sequência de citações; os trabalhos citados devem ser interpretados, analisados e, se possível, comparados com o trabalho em questão, demonstrando a conveniência da confirmação ou continuação do tema.

Esse capítulo também pode incluir uma revisão da literatura que busca descrever trabalhos desenvolvidos por outros atores sobre o mesmo assunto do TCC ou assunto similar. A análise desta literatura permite realizar comparações entre o tema escolhido e os trabalhos anteriores, evitando-se repetições desnecessárias e registrando-se o estado da arte do tema selecionado. Ao citar a literatura correlata ao assunto, reforçam-se a necessidade e a oportunidade do estudo.

Nesse sentido, a literatura revisada deve ser interpretada e confrontada com os objetivos deste trabalho, e não apenas listada.
Essa análise crítica permite identificar o estado da arte do tema, justificar a pesquisa e evitar repetições.
Cabe ressaltar, ainda, a necessidade de rigor no registro completo das fontes, incluindo o DOI dos artigos e as URLs das páginas acessadas.
Exemplo de citação com um trabaho produzido no Ibmec~\cite{assumpcao2024estacionamento}